\emph{Conventions.} $\Cl(1,3)$ with signature $(+,-,-,-)$; vectors $p = p^{\mu}\gamma_{\mu}$, $\gamma^{\mu}$ the reciprocal frame; the trace of the Dirac matrices is $4\langle\cdot\rangle_{0}$ and the units of the Feynman rules are coefficients in front of the product (\claimref{synthesis.one-loop-qed-in-the-spacetime-algebra}(1), (5)). The vertex is $\gamma^{\mu}$ with the coupling of \claimref{field-theory.coupled-system-tree-level}; the free Dirac fields and their two spin states per momentum are those of \claimref{field-theory.second-quantisation}. The s-channel photon at tree level couples conserved currents, so the stored photon suffices there. The polarisation loop inside the photon line is a fermion loop with no internal photon; the stored result computes it with Pauli-Villars fermions (Fact \claimref{synthesis.imported-one-loop-integrals-and-pauli-villars-regulators}), and the subtracted function $\hat\Pi$ used here does not depend on the regulator (a regulator-free construction with no Krein form is \claimref{synthesis.one-loop-vacuum-polarisation-extends-without-a-regulator-and-ward-preserving-extensions-differ-only-by-the-renormalisation-constants}(ii)). Only the charge-renormalisation sentence of Step 5 uses loops with an internal photon, and so the Krein-form premise \claimref{synthesis.krein-form-premise-indefinite-inner-product-for-gauge-fields}.

\emph{Step 1 (LSZ).} Let $Q(p) = m^{2}$. Then $p^{2} = m^{2}$ as multivectors, so $(p + m)^{2} = p^{2} + 2mp + m^{2} = 2m(p + m)$, $(p - m)^{2} = -2m(p - m)$ and $(p + m)(p - m) = p^{2} - m^{2} = 0$; $\langle p\rangle_{0} = 0$ gives $4\langle p \pm m\rangle_{0} = \pm 4m$. So for $m \neq 0$, $P_{\pm} = (m \pm p)/(2m)$ are complementary idempotents with $4\langle P_{\pm}\rangle_{0} = 2$: each projects the spinor module onto a two-dimensional space, the two spin states. Off shell $(p - m)S(p) = 1$ with $S(p) = (p + m)/(Q(p) - m^{2})$, and $(Q(p) - m^{2})S(p) = p + m$, the residue at the pole. The reduction formula \cite{LSZ1955} (Fact \claimref{synthesis.imported-lsz-reduction-ward-identity-and-the-second-qed-beta-coefficient}(i); Fact \claimref{synthesis.imported-one-loop-integrals-and-pauli-villars-regulators}(v)) writes an S-matrix element as the amputated Green function at the poles times $Z_{2}^{1/2}$ per external fermion; at tree level $Z_{2} = 1$. Summed over spins, an outgoing particle and an incoming particle contribute $p + m$ (the residue), an incoming antiparticle and an outgoing antiparticle $p - m$ (the residue of $S(-p)$ up to sign). For $m = 0$ the residue $p$ is nilpotent ($p^{2} = 0$) and $4\langle\cdot\rangle_{0}$ of $p$ times anything is the massless spin sum. The check verifies these identities at three rational on-shell momenta (masses $4$, $12$, $0$) and the residue at an off-shell momentum.

\emph{Step 2 (tree amplitude).} With $q = p_{1} + p_{2}$, the amplitude is $M = \frac{e^{2}}{s}\,[\bar v(p_{2})\gamma^{\mu}u(p_{1})][\bar u(p_{3})\gamma_{\mu}v(p_{4})]$ (units collected in front). Squaring and summing over spins by Step 1, $\tfrac14\sum|M|^{2} = \tfrac{e^{4}}{4s^{2}}T_{e}^{\mu\nu}T_{\mu\nu}^{\mu}$ with $T_{e}^{\mu\nu} = 4\langle(p_{2} - m_{e})\gamma^{\mu}(p_{1} + m_{e})\gamma^{\nu}\rangle_{0}$, $T^{\mu}_{\mu\nu} = 4\langle(p_{3} + m_{\mu})\gamma_{\mu}(p_{4} - m_{\mu})\gamma_{\nu}\rangle_{0}$. Expanding $4\langle abcd\rangle_{0} = 4[(a\cdot b)(c\cdot d) - (a\cdot c)(b\cdot d) + (a\cdot d)(b\cdot c)]$ and $4\langle ab\rangle_{0} = 4a\cdot b$ gives the standard form of clause (2); the check computes the two traces and their contraction exactly at three rational configurations (both masses, electron massless, both massless) and compares with the standard form, with a control (the particle residue in place of the positron's fails). Massless kinematics in the centre-of-mass frame, $p_{1} = E(\gamma_{0} + \gamma_{3})$, $p_{3} = E(\gamma_{0} + \sin\theta\,\gamma_{1} + \cos\theta\,\gamma_{3})$, give $\tfrac14\sum|M|^{2} = e^{4}(1 + \cos^{2}\theta)$ (checked at five rational angles; a polynomial of degree two in $\cos\theta$ is fixed by three). With $d\sigma/d\Omega = \tfrac{1}{64\pi^{2}s}\tfrac{|p_{3}|}{|p_{1}|}\tfrac14\sum|M|^{2}$ (standard two-body phase space \cite{PeskinSchroeder1995}) and $e^{2} = 4\pi\alpha$: $d\sigma/d\Omega = \tfrac{\alpha^{2}}{4s}(1 + \cos^{2}\theta)$ and $\sigma = 2\pi\tfrac{\alpha^{2}}{4s}\int_{-1}^{1}(1 + c^{2})\,dc = \tfrac{4\pi\alpha^{2}}{3s}$. For $m_{e} = 0$ and muon speed $\beta$ ($m_{\mu} = \tfrac{\sqrt s}{2}(1 - \beta^{2})^{1/2}$) the trace result is a polynomial of degree two in $c = \cos\theta$; it is fixed by $c = -1, 0, 1$, checked at $c = 3/5$, and $\sigma\,s/(\pi\alpha^{2}) = \tfrac{\beta}{2}\int_{-1}^{1}\tfrac{|M|^{2}}{e^{4}}dc = \tfrac43\cdot\tfrac{\beta(3 - \beta^{2})}{2}$ at $\beta = \tfrac35, \tfrac45, \tfrac{12}{13}, 1$, exactly. Since both sides are polynomials in $\beta$ of degree three after multiplying by $s^{2}$ and fixing $s$, four values of $\beta$ establish the identity at that $s$, and both scale as $1/s$.

\emph{Step 3 (one loop: the insertion).} The photon two-point function with polarisation insertions is, in a covariant gauge, $\frac{-i}{q^{2}(1 - \hat\Pi(q^{2}))}\bigl(g_{\mu\nu} - \frac{q_{\mu}q_{\nu}}{q^{2}}\bigr) + (\text{gauge part})\frac{q_{\mu}q_{\nu}}{q^{4}}$, by transversality of the stored polarisation (\claimref{synthesis.one-loop-qed-in-the-spacetime-algebra}(2)). Between on-shell currents every $q_{\mu}$ term vanishes: $4\langle(p_{2} - m)\,q\,(p_{1} + m)\gamma^{\nu}\rangle_{0} = 0$ for $q = p_{1} + p_{2}$, because $(p_{2} - m)(p_{1} + p_{2})(p_{1} + m) = (p_{2} - m)[(p_{1} - m) + (p_{2} + m)](p_{1} + m) = 0$ by Step 1; the check verifies this for both currents and every $\nu$, with the control $q = p_{1} - p_{2}$, which does not drop. So the s-channel amplitude is $M_{\mathrm{tree}}/(1 - \hat\Pi(s))$, with $\hat\Pi(s) = \Pi(s) - \Pi(0)$ the on-shell subtracted function of the stored result: $\hat\Pi(s) = -\tfrac{2\alpha}{\pi}\int_{0}^{1}x(1 - x)\ln\frac{m^{2}}{m^{2} - x(1-x)s - i0}dx$ (\cite{PeskinSchroeder1995}, sec. 7.5, eqs. (7.90) to (7.92); the regulator scale cancels in the subtraction).

\emph{Step 4 (absorptive part and the optical theorem).} For $s > 4m^{2}$ the argument $m^{2} - x(1-x)s$ is negative exactly on the window $x_{-} < x < x_{+}$, $x_{\pm} = (1 \pm \beta)/2$, since $x_{\pm}(1 - x_{\pm}) = (1 - \beta^{2})/4 = m^{2}/s$ (checked). There $\ln\frac{m^{2}}{A - i0} = \ln\frac{m^{2}}{|A|} + i\pi$, so $\mathrm{Im}\,\hat\Pi(s + i0) = -2\alpha\int_{x_{-}}^{x_{+}}(x - x^{2})dx = -\tfrac{\alpha}{6}\beta(3 - \beta^{2})$, exactly (checked at four rational $\beta$; the integral is the polynomial $\beta(3 - \beta^{2})/12$). By Step 2, $s\,\sigma_{\mathrm{tree}}(e^{+}e^{-} \to f\bar f)/(4\pi\alpha) = \tfrac{\alpha}{4}\cdot\tfrac43\cdot\tfrac{\beta(3 - \beta^{2})}{2} = \tfrac{\alpha}{6}\beta(3 - \beta^{2})$, the same. This is the s-channel optical-theorem relation for the photon self-energy: the discontinuity of the one-loop photon line is the phase-space integral of the squared tree amplitude for producing the pair (\cite{PeskinSchroeder1995}, sec. 7.3; argued here, the two sides computed). Full unitarity of the S-matrix at order $e^{4}$ also involves $e^{+}e^{-} \to \gamma\gamma$, the t-channel and box cuts, which are not computed.

\emph{Step 5 (large $s$ and the charge).} For $s \gg m^{2}$, $\ln\frac{m^{2}}{|m^{2} - x(1-x)s|} = -\ln\frac{s}{m^{2}} - \ln(x(1-x)) + O(m^{2}/s)$, so $\mathrm{Re}\,\hat\Pi = \tfrac{2\alpha}{\pi}[\tfrac16\ln\tfrac{s}{m^{2}} + J]$ with $J = \int_{0}^{1}x(1-x)\ln(x(1-x))\,dx = 2\int_{0}^{1}(x - x^{2})\ln x\,dx = 2(-\tfrac14 + \tfrac19) = -\tfrac{5}{18}$, by $\int_{0}^{1}x^{n}\ln x\,dx = -1/(n+1)^{2}$: $\mathrm{Re}\,\hat\Pi = \tfrac{\alpha}{3\pi}(\ln\tfrac{s}{m^{2}} - \tfrac53)$ (the rationals checked). In multiplicative renormalisation the charge is $e = Z_{1}^{-1}Z_{2}Z_{3}^{1/2}e_{0}$ (the vertex counterterm $Z_{1}$, the fermion field factor $Z_{2}^{1/2}$ on either side of the vertex, which LSZ supplies on external legs, and the photon field factor $Z_{3}^{1/2}$); under the Krein-form premise, with $Z_{1} = Z_{2}$ (\claimref{synthesis.one-loop-qed-in-the-spacetime-algebra}(4); Ward's identity, Fact \claimref{synthesis.imported-lsz-reduction-ward-identity-and-the-second-qed-beta-coefficient}(ii)) the fermion factors cancel and $e = Z_{3}^{1/2}e_{0}$, so the insertion is the running coupling $\alpha(s) = \alpha/(1 - \mathrm{Re}\,\hat\Pi(s))$ at this order (argued). Not computed: the vertex and box corrections and real soft-photon emission, whose infrared divergences cancel only in the inclusive rate.

\emph{Status.} Clause (1) proved (algebra); clause (2) computed exactly (an identity of polynomials, sampled at rational kinematics) with the phase-space formula imported; clause (3) computed for the absorptive part and the constant, argued for the optical-theorem reading and the charge renormalisation. All standard (LSZ 1955; Peskin and Schroeder secs. 5.1, 7.3, 7.5), recomputed in the spacetime algebra. Transport-absent; the spacetime algebra is convenient.
